Recent large language models (LLMs) have demonstrated strong capabilities in
understanding and generating code~\cite{chen2021evaluating,austin2021program,nijkamp2022codegen},
achieving strong performance in tasks ranging from competitive
programming~\cite{li2022competition} to repository-level software
engineering~\cite{jimenez2023swe}.
Building on these capabilities, the role of code in agentic systems is
expanding beyond a target artifact to be generated.
Programs are increasingly used as the medium through which LLM
agents reason, act, and model their environments.
Program-aided reasoning methods externalize intermediate
computation into executable code~\cite{chen2022program,gao2023pal,li2023chain};
robotic and embodied agents use generated programs as executable
policies for interacting with physical or simulated
worlds~\cite{ahn2022can,liang2023code};
and software-engineering or interactive environments use
codebases, execution traces, tests, and runtime feedback as
structured representations of environment state and dynamics, in
which agents plan, act, and revise their behavior~\cite{yang2023intercode,jimenez2023swe,liu2023agentbench}.
Taken together, these developments suggest a broader view:
code is not only an artifact generated by LLMs, but also an
executable, inspectable, and stateful medium through which agents
reason, act, observe feedback, and verify progress. We refer to
this view as \emph{code as agent harness}.



\begin{figure}[t]
    \centering
    \includegraphics[width=1.0\linewidth]{figures/overview.pdf}
    \caption{Taxonomy of code as agent harness.}
    \label{fig:taxonomy}
    \vspace{-5mm}
\end{figure}

Recent discussions on \emph{agent harnesses}~\cite{lee2026metaharness,lou2026autoharness,anthropic2025longrunning,lopopolo2026harnessengineering}
provide a useful system-level lens for understanding this shift.
An agent harness refers to the software layer that surrounds an LLM
with tools, APIs, sandboxes, memory, validators, permission
boundaries, execution loops, and feedback channels, thereby turning
a stateless model into a functional agent capable of long-running
task execution~\cite{zhang2025agentic,agrawal2025gepa,zhang2023toolcoder,wang2025teaching,lavon2025execution,cheng2026llm,dai2025feedbackeval}.
In this view, the bottleneck of autonomy is not only the reasoning
ability of the base model, but also the reliability of the system
that connects model outputs to long-horizon actions and persistent states.

To clarify the role of code in this broader harness view, we
distinguish three coupled elements of long-running agentic systems:
\emph{model-internal capabilities}, \emph{system-provided harness
infrastructure}, and \emph{agent-initiated code artifacts}.
\emph{Model-internal capabilities} refer to the model's reasoning,
perception, planning, simulation, and evaluation abilities.
\emph{System-provided harness infrastructure} refers to the
predefined tools, APIs, sandboxes, memory systems, validators,
permission boundaries, telemetry, and workflows that connect model
outputs to external actions and feedback, and forms the main focus
of harness engineering~\cite{openai2026harnessengineering,langchainanatomyharness2026}.
In contrast, \emph{agent-initiated code artifacts}, which remain
relatively underexplored, are interactive code objects that agents
create, execute, observe, revise, persist, and share within the task
execution loop. Through execution feedback, these artifacts help
agents reason, act, verify progress, store state, and coordinate
with other agents. Examples include regression tests, temporary
tools, DSL programs, executable workflows, reusable skills, and
intermediate program states. Representative systems such as Claude
Code~\cite{claudecode2025}, Codex~\cite{codex2025}, LangChain
~\cite{langchaindeepagentsharness2026}, and enterprise agent
platforms show how these elements jointly enable adaptation in long-running
agent systems.



With this distinction in mind, we revisit the role of code
in agentic systems. Existing surveys typically either treat code as the end product of
LLMs. In contrast, we focus on \emph{agent-initiated code artifacts} and
how model capabilities construct and evolve them through
interaction with harness infrastructure, with code serving as the
organizing center for the interface, agent capabilities,
and multi-agent coordination.
Across diverse agentic systems, code is used not only to produce
solutions, but also to execute reasoning, ground actions, maintain
state, and expose feedback. We term this view \emph{code as agent
harness}: code as the executable and inspectable medium through
which agents reason, act, and adapt. This shifts the scope from producing correct programs to
understanding how code supports reliable closed-loop agentic
behavior.



To systematically characterize \emph{code as agent harness}, we
organize the survey into three connected layers, as shown in
Figure~\ref{fig:taxonomy}.
This organization follows how code becomes an operational medium
inside the agent loop: it first enters as a harness interface
for reasoning, acting, and environment representation; it then
supports harness mechanisms that manage planning, memory, tool
use, execution, and repair over time; and it finally becomes a
shared artifact through which multiple agents coordinate over
repositories, tests, traces, workflows, and execution states.

First, \emph{\textbf{Harness Interface: Code for Reasoning, Acting, and
Environment Modeling}} (\S\ref{sec:foundations}) studies how code
forms the basic interface between a model and its task
environment.
At this layer, code is the medium that converts model outputs into
executable and inspectable structures.
We review \emph{code for reasoning}, where programs externalize
intermediate computation and allow interpreters, symbolic solvers,
execution traces, or process rewards to check and refine reasoning
~\cite{gao2023pal,chen2022program,li2023chain,ye2023satlm,ni2024next,li2025codeprm}.
We then review \emph{code for acting}, where generated programs
serve as policies, tool calls, behavior trees, or reusable skills
for embodied, GUI, and software environments
~\cite{ahn2022can,liang2023code,wang2023voyager,mu2024robocodex,zhang2025codebt,lin2026ui}.
Finally, we examine \emph{code for environment modeling}, where
program states, repositories, traces, simulators, and tests
represent state, dynamics, and feedback signals for agent
interaction
~\cite{tang2024worldcoder,copet2025cwm,zheng2026code2world,jimenez2023swe,liu2023agentbench,gandhi2026endless}.
This layer establishes the core harness interface: code is how the
agent makes reasoning executable, action programmable, and
environment state inspectable.

Building on this interface, \emph{\textbf{Harness Mechanisms:
Planning, Memory, Tool Use, Control, and Optimization}}~(\S\ref{sec:modules})
studies how code-harnessed agents remain reliable beyond a single
generation step.
Once code is placed inside the agent loop, the harness must decide
what to execute next, preserve useful state, expose the right
tools, and convert failures into corrective actions.
We therefore review planning methods that organize long-horizon
software tasks through decomposition, structural grounding,
trajectory search, or workflow orchestration
~\cite{jiang2024selfplanning,gur2023webagent,bairi2024codeplan,li2025codetree,islam2024mapcoder};
memory methods that maintain working state, retrieve repository
evidence, store reusable experience, and support shared
interaction histories
~\cite{gaurav2025codemem,zhang2024autocoderover,zhang2023repocoder,wang2026memgovern};
tool-use methods that connect agents to APIs, repositories,
execution environments, and verification tools
~\cite{zhang2023toolcoder,liu2024toolnet};
and feedback-driven control and harness optimization methods that use static analysis,
runtime errors, tests, and human feedback to revise code through
repeated execution
~\cite{huang2023agentcoder,ukai2024adacoder,Nunez2024AutoSafeCoder,li2026agentharness}.
This layer turns the interface in \S\ref{sec:foundations} into an
operational harness: planning controls the execution trajectory,
memory preserves state, tools expand the action space, and
feedback-driven adaptation closes the loop between failure and revision.

Finally, \emph{\textbf{Scaling the Harness: Multi-Agent Orchestration over Code}} (\S\ref{sec:mas}) extends the harness from
a single agent to collaborative ecosystems.
When multiple agents operate over code, the harness must not only
support individual reasoning and execution, but also coordinate
roles, share intermediate artifacts, maintain common state, and
verify collective progress.
We review multi-agent code-centric systems through agent roles
such as manager, planner, coder, reviewer, and tester;
collaboration modes such as programming, repair, debate,
red-teaming, and adversarial interaction; and workflow topologies
ranging from centralized coordination to distributed or streaming
collaboration
~\cite{wu2024autogen,Hong2023MetaGPT,Dong2024SelfCollaboration}.
This layer shows how code becomes a shared harness for
orchestrated autonomy: repositories, tests, traces, and structured
artifacts provide the common workspace through which agents
coordinate, inspect, and improve each other's behavior.

\begin{tcolorbox}[
  agentscope,
  float,
  floatplacement=t
]
This survey studies \emph{code as agent harness}: code-centered agent systems where reasoning, action, state, feedback,
and verification are organized around executable, inspectable, and
stateful programs.
We organize the literature up to 2026 into three connected layers:
\begin{itemize}
    \item \textbf{Harness Interface}: code enters the agent loop
    as a reasoning substrate, an action interface, and an
    environment representation.

    \item \textbf{Harness Mechanisms}: planning, memory, tool use,
    control, and harness optimization sustain code-centric agents over
    long-horizon execution and revision.

    \item \textbf{Scaling the Harness}: shared code artifacts,
    execution states, repositories, and structured workflows
    support coordination, review, and collective verification in
    multi-agent systems.
\end{itemize}


\end{tcolorbox}

Beyond the taxonomy, we examine how agent-initiated code interaction appears across five application domains. In coding assistance, agents author patches, tests, and issue-resolution workflows over live repositories~\cite{jimenez2023swe,yang2024swe,wang2024openhands}. In GUI and OS automation, agents synthesize and execute interface commands grounded in DOM trees, accessibility APIs, and executable evaluators~\cite{deng2023mind2webgeneralistagentweb,zhou2024webarenarealisticwebenvironment}. In scientific discovery, agents dynamically compose and execute hypothesis-testing pipelines spanning simulations, lab protocols, and data analysis~\cite{bran2023chemcrowaugmentinglargelanguagemodels,boiko2023autonomous,lu2024aiscientistfullyautomated,huang2025biomni}. In personalization and embodied control, agents author and revise executable policies, simulators, and skill libraries in response to environment feedback~\cite{ahn2022can,liang2023code,wang2023voyager}. We further outline open challenges for harness engineering, including evaluation beyond final task success, verification under incomplete feedback, regression-free harness improvement, consistent shared state across multiple agents, human oversight, and extensions to multimodal environments. This survey provides a roadmap for studying code not only as something agents generate, but as the runtime medium through which they execute, adapt, and coordinate reliable behavior.


\begin{tcolorbox}[agentcontrib]
% This survey makes the following contributions:
\begin{itemize}
    \item \textbf{Conceptual framing}: We formalize
    \emph{code as agent harness}, reframing code from a generated
    artifact into the operational substrate of executable,
    verifiable, and stateful AI agent systems.

     \item \textbf{Taxonomy and synthesis}: We organize
   code as agent harness into three connected layers: harness
    interfaces, harness mechanisms, and scaling harness, and synthesize representative methods.

    \item \textbf{Applications and future agenda}:We connect the taxonomy
to real-world applications and outline key challenges in evaluation,
verification, safety, and coordination.
\end{itemize}
\end{tcolorbox}
